\documentclass[10pt,a4paper]{amsart}
\usepackage[margin=19mm]{geometry}
\usepackage{amsmath,amssymb,mathtools}
\usepackage[scr=rsfs,cal=euler]{mathalfa}
\usepackage{xfrac}
\usepackage[unicode,hidelinks,bookmarks=false]{hyperref}
\newcommand{\ie}{i.e.\ }
\newcommand{\cf}{cf.\ }
\newcommand{\resp}{respectively\ }
\newcommand{\Subsection}{\subsection}
\providecommand{\nobreakdash}{\nobreak\hbox{-}}
\setlength{\emergencystretch}{2em}
\multlinegap0pt
\usepackage{xcolor}
\usepackage{amssymb}
\usepackage{tikz}
\usepackage[normalem]{ulem}
\newtheorem{lem}{Lemma}
\newtheorem{thm}{Theorem}
\usepackage{cleveref}
\crefname{lem}{Lemma}{Lemmas}
\crefname{thm}{Theorem}{Theorems}
\def\cC{\mathcal{C}}
\def\cD{\mathcal{D}}
\DeclareMathOperator{\tk}{Tk}
\title{Adjoints of Morphisms of Neural Codes}
\author{Juliann Geraci, Alexander B. Kunin, Alexandra Seceleanu}
\date{Source: arXiv:2603.10837v1 (CC BY 4.0)}
\begin{document}
\maketitle

\noindent\emph{Source and licence.} Adjoints of Morphisms of Neural Codes. Version arXiv:2603.10837v1 (CC BY 4.0), distributed by arXiv under the Creative Commons Attribution 4.0 International licence, http://creativecommons.org/licenses/by/4.0/, as recorded in the arXiv licence metadata for this exact version and shown on the abstract page https://arxiv.org/abs/2603.10837v1. Full licence notice: \texttt{licenses/arxiv-2603.10837-CC-BY-4.0.txt}.\par\smallskip

\noindent\textit{Editorial scope.} Counted unit: the article's thm lem:bmfredundancy (source0.tex lines 733--737). The article's complete original proof of it is delivered with it (source0.tex lines 738--771, one \texttt{proof} environment of 34 lines). No mathematics is written, repaired or re-derived: the statement and the proof are the article's own lines, in the article's own notation, and no retained line is substituted or re-flowed.  Retained uncounted as the source-local context the proof needs: lines 724--728; lines 1138--1140.  Nothing was omitted from any retained span, so the package carries no omission record.  No retained line contains a citation, so the wrapper carries no bibliography and no bibliography entry.  No retained line contains a figure environment, an image-inclusion command or drawing code, so no image is delivered and the package declares no asset dependency.  Editorial formatting only, in addition: an AMS \texttt{amsart} preamble that restates the article's own declarations and macros used by the retained lines, namely the environments \texttt{lem}, \texttt{thm} on separate counters, together with the standard packages those lines need.  The retained environments print in this document as Lemma 1, Theorem 1 in order of appearance, whereas the article prints the same items with the numbering of its own full document.  The complete native source of the article as distributed in the arXiv e-print for this exact version is bundled unchanged as \texttt{sources/196171/source0.tex}.
        \begin{lem}\label{lem:redundancyinimage}
            Let $f: \cC \subseteq2^{[n]} \to \cD\subseteq2^{[k]}$ be a surjective morphism of codes, and suppose $k$ is redundant in $\cD$.
            For any representative $t = (\tau_1,\dots,\tau_k)$ of $f$, the tuple $t' = (\tau_1,\dots,\tau_{k-1})$ represents composition of $f$ with deletion of neuron $k$ in $\cD$.
            In particular, the image of $f_{t'}$ is isomorphic to the image of $f$.
        \end{lem}

    \begin{lem}\label{lem:rootsofintersections}
        Let $S,T$ be two trunks of a code $\cC$. Then $\sqrt{S \cap T} \supseteq \sqrt{S} \cup \sqrt{T} \supseteq \sqrt{S} \cap \sqrt{T}$.
    \end{lem}

        \begin{thm}\label{lem:bmfredundancy}
            Suppose $f: \cC \to \cD$ is a BMF with $\cC \subseteq 2^{[n]}$, $\cD \subseteq 2^{[k]}$.
            If $\cC$ is reduced, 
            then there is a BMF $\cC \to \rho(\cD)$. % where $\cD'$ is reduced and $\cD' \cong \cD$.
        \end{thm}

        \begin{proof}
            We will show that, if $k$ is redundant to $\sigma \subseteq[k-1]$ in $\cD$, then composing $f$ with the deletion of neuron $k$ is a BMF.

            For each $i \in [k]$, set $T_i = f^{-1}(\tk_\cD(i))$, and set $t = (\sqrt{T_1},\dots,\sqrt{T_k})$ a representative of $f$. Our goal is to show that the map represented by $t' = (\sqrt{T_1},\dots,\sqrt{T_{k-1}})$ is also a Boolean matrix factorization. Since $k$ is redundant in $\cD$, from \cref{lem:redundancyinimage} we know that the image of  $f_{t'}$ is exactly $\cD$ with $k$ deleted, so all that remains to show is that $f_{t'}^\top (f_{t'}(c)) = c$ for all $c \in \cC$.
            
            First, we claim that $f_{t'}^\top(f_{t'}(c)) = f^\top(f(c)\setminus k)$.
            That $f_{t'}(c) = f(c) \setminus k$ follows from the definition of $f_{t'}$.
            Writing $f_{t'}^\top(d) = \bigcup_{j\in d} \sqrt{T_j}$ we see that this agrees with $f^\top(d)$ when $k \notin d$.
            
            Consider some $c \in \cC$ and let $d = f(c)$.
            If $c \notin T_k$, then $k \notin d$, so from above we have $f_{t'}(c) = f(c)$ and $f_{t'}^\top (f_{t'}(c)) = c$.

            Now consider the case $c \in T_k$, so $k \in d$.
            Since $k$ is redundant to $\sigma$, $d \in \tk_\cD(\sigma)$ and so we can write $d$ as the disjoint union $d = d' \cup \sigma \cup k$. Then
            \begin{align*}
                f^\top(d) = \bigcup_{j\in d'} \sqrt{T_j} \cup \bigcup_{j\in\sigma} \sqrt{T_j} \cup \sqrt{T_k}\tag{*}\label{eq:needtoreducefactor}
            \end{align*}
            We will show that $\sqrt{T_k}$ is unnecessary in this union.
            
            Since $k$ is redundant to $\sigma$, we have $\tk_\cD(k) = \tk_\cD(\sigma) = \bigcap_{j\in\sigma} \tk_\cD(j)$.
            Applying $f^{-1}$ to both sides we have $T_k = \bigcap_{j\in\sigma} T_j$ and so, taking the root of both sides and applying \cref{lem:rootsofintersections},
            \[ \sqrt{T_k} = \sqrt{\bigcap_{j\in\sigma} T_j} \supseteq \bigcup_{j\in\sigma} \sqrt{T_j}. \]
            If $\sqrt{T_k} = \bigcup_{j\in\sigma} \sqrt{T_j}$, then clearly $\sqrt{T_k}$ is unnecessary in \eqref{eq:needtoreducefactor}.

            Suppose, on the other hand, there is some $\ell \in \sqrt{T_k} \setminus \bigcup_{j\in\sigma}\sqrt{T_j}$. If $\tk_\cC(\ell) = \tk_\cC(\sqrt{T_k}) =T_k$
            then $\ell$ is redundant to $\bigcup_{j\in\sigma} \sqrt{T_j}$
            because
                \[ \tk_\cC\Bigl(\bigcup_{j\in\sigma} \sqrt{T_j}\Bigr) = \bigcap_{j\in\sigma} \tk_\cC(\sqrt{T_j}) = \bigcap_{j\in\sigma} T_j = T_k = \tk_\cC(\ell). \]
            Since $\cC$ is reduced, therefore, we must have $\tk_\cC(\ell) \supsetneq T_k$.
            Consider $\alpha \in \tk_\cC(\ell) \setminus T_k$. Since $\ell \in \alpha$, $k \notin f(\alpha)$, and $f^\top(f(\alpha)) = \alpha$, we see from \eqref{eq:needtoreducefactor} that there must be some $j \in d'$ such that $\ell \in \sqrt{T_j}$. Once again, we conclude that $\sqrt{T_k}$ is unnecessary in \eqref{eq:needtoreducefactor}.

            Therefore, $f_{t'}$ is a BMF.
            The statement of the lemma follows by permuting neurons in $\cD$ if necessary and inducting on the number of redundant neurons in $\cD$.
        \end{proof}

\end{document}
